% Bagian: first-order-logic
% Bab: syntax
% Subbagian: intro-syntax

\documentclass[../../../include/open-logic-section]{subfiles}

\begin{document}

\olfileid[id]{fol}{syn}{itx}

\olsection{Pendahuluan}

Untuk mengembangkan teori dan metateori logika orde pertama, mula-mula kita
harus mendefinisikan sintaksis dan semantik ungkapannya. Ungkapan logika orde
pertama berupa suku dan !!{formula}s. Suku dibentuk dari !!{variable}s,
!!{constant}s, dan !!{function}s. Selanjutnya, !!^{formula}s dibentuk dari
!!{predicate}s bersama suku (yang membentuk !!{formula}s terkecil, yakni
``atomik''), lalu dari !!{formula}s atomik kita dapat membentuk formula yang
lebih kompleks dengan menggunakan penghubung logis dan kuantor. Ada banyak
cara berbeda untuk menetapkan aturan pembentukan; di sini kita hanya
memberikan salah satu cara yang mungkin. Sistem lain akan memilih simbol
yang berbeda, memilih himpunan penghubung primitif yang berbeda, dan
menggunakan tanda kurung secara berbeda (atau bahkan tidak menggunakannya
sama sekali, seperti dalam apa yang disebut notasi Polandia). Meskipun
demikian, semua pendekatan memiliki kesamaan: aturan pembentukan
mendefinisikan himpunan suku dan !!{formula}s secara \emph{induktif}. Jika
dilakukan dengan tepat, setiap ungkapan pada dasarnya hanya dapat dihasilkan
melalui satu cara menurut aturan pembentukan. Definisi induktif yang
menghasilkan ungkapan yang \emph{terbaca secara unik} memungkinkan kita
memberikan makna kepada ungkapan tersebut dengan metode yang sama---definisi
induktif.

\end{document}

